\documentclass[12pt]{article}
\usepackage[T1]{fontenc}
\usepackage[utf8]{inputenc}
\usepackage{lmodern}
\usepackage{textcomp}
\usepackage{enumitem}
\usepackage{geometry}
\usepackage{graphicx}
\usepackage{amsmath, amssymb}
\geometry{margin=1in}
\begin{document}
\begin{center}
History - K-12 Level Assessment 17 \
Total Marks: 40 \
Answer all questions.
\end{center}

\section*{Multiple Choice (10 marks)}
\begin{enumerate}[label=\arabic*.]
\item The early Mesolithic Maglemosian culture was adapted to:
\begin{enumerate}[label=(\alph*)]
    \item a tundra environment.
    \item a desert environment.
    \item a forest and lakeside environment.
    \item an equatorial environment.
\end{enumerate}
\item Researchers who study humans by residing in particular societies and observing the behaviors of the people are known as:
\begin{enumerate}[label=(\alph*)]
    \item archaeologists.
    \item ethnographers.
    \item linguists.
    \item paleoanthropologists.
\end{enumerate}
\item Excavation of Yang-shao sites in China indicates which domesticated crops?
\begin{enumerate}[label=(\alph*)]
    \item beans, millet, and maize.
    \item wheat, rice, and peas.
    \item sorghum, emmer, and legumes.
    \item millet, cabbage, and rice.
\end{enumerate}
\item Which of the following describes a key change in hominids beginning at least as early as Homo erectus that is probably related to increasingly larger brain size?
\begin{enumerate}[label=(\alph*)]
    \item microcephaly
    \item neoteny
    \item prognathism
    \item supraorbital cortex
\end{enumerate}
\item Of the following, which is the earliest example of the development of a complex society?
\begin{enumerate}[label=(\alph*)]
    \item Jericho in Israel
    \item Caral in Peru
    \item Stonehenge in England
    \item Olmec along the Gulf Coast
\end{enumerate}
\end{enumerate}

\section*{True / False (10 marks)}
\textit{Answer True or False}
\begin{enumerate}[label=\arabic*.]
\setcounter{enumi}{5}
\item True or False: The recovery of over 6,500 bones from a single hominid species at the Cave of the Bones provides unusually good evidence of diversity.
\item True or False: The Cordilleran ice mass was a significant Pleistocene glacier centered in the Rocky Mountains of North America.
\item True or False: Since the late 1800 s, mean global sea level has increased by about 20 centimeters (8 inches) and will continue to rise.
\item True or False: A landfill where refuse from an entire community has accumulated represents a primary context.
\item True or False: The Inca developed one of the earliest kingdoms in South America, reaching its height around AD 200.
\end{enumerate}

\section*{Long Form Response (20 marks)}
\textit{Answer each question with a clear and complete explanation.}
\begin{enumerate}[label=\arabic*.]
\setcounter{enumi}{10}
\item Discuss the evidence that supports the idea that Neandertals were hunters, focusing on the significance of animal bones with stone tool cut marks and the absence of animal teeth marks.
\item Discuss the reasons behind the Mogollon's decision to leave their mountain homeland, focusing on the impact of environmental factors such as drought on their way of life.
\end{enumerate}

\vfill
\noindent\textit{End of Paper}
\end{document}